\subsection{Horizontální pohyb a akcelerace}

Metoda \texttt{Run()} aplikuje horizontální pohyb pomocí \texttt{Lerp} interpolace, která zajišťuje plynulé zrychlení a zpomalování namísto okamžité změny rychlosti. Vzorec \texttt{Mathf.Lerp(rb.linearVelocity.x, targetSpeed, Time.fixedDeltaTime * accelRate)} postupně přesouvá aktuální rychlost k cílové rychlosti.

Zrychlení se liší podle toho, zda je hráč na zemi nebo ve vzduchu. Parametr \texttt{accelInAir} určuje, kolik procent pozemního zrychlení se použije ve vzduchu. Typicky se používá nižší hodnota (například 50-70 procent), což zajišťuje, že hráč nemůže ve vzduchu měnit směr tak rychle jako na zemi. Toto omezení přidává hloubku do herní mechaniky a odměňuje hráče za správné načasování pohybu.

Během provádění wall jump je volání \texttt{Run()} dočasně deaktivováno, aby hráč nemohl během letu po stěně reagovat na směrové tlačítko. Toto zajišťuje, že wall jump má konzistentní trajektorii bez ohledu na input hráče během letu.

\begin{lstlisting}[language=csharp,caption={Listing 4: Metoda Run pro horizontální pohyb}]
private void Run()
{
    if (pw != null && pw.IsWallJumping)
        return;
    if (rb == null || data == null) return;

    float targetSpeed = moveInput.x * data.runMaxSpeed;

    float accelRate = lastOnGroundTime > 0 ?
        data.runAccelAmount :
        data.runAccelAmount * data.accelInAir;

    if (isJumping && Mathf.Abs(rb.linearVelocity.y) < data.jumpHangTimeThreshold)
    {
        accelRate *= data.jumpHangAccelerationMult;
        targetSpeed *= data.jumpHangMaxSpeedMult;
    }

    float newVelX = Mathf.Lerp(rb.linearVelocity.x, targetSpeed, 
        Time.fixedDeltaTime * accelRate);
    rb.linearVelocity = new Vector2(newVelX, rb.linearVelocity.y);
}
\end{lstlisting}
